\chapter{Annexes}
\label{chap:annexes}

\section{Guide d'installation}
\label{sec:guide_installation}

\subsection{Prérequis}
\begin{itemize}
    \item Node.js v18+ et npm
    \item PostgreSQL v14+
    \item Python 3.10+ avec pip
    \item Flutter SDK 3.x
    \item Git
\end{itemize}

\subsection{Installation du backend}
\begin{lstlisting}[language=bash,caption={Installation du backend}]
cd backend
npm install
cp .env.example .env          # Configurer les variables
npx prisma generate           # Générer le client Prisma
npx prisma migrate deploy     # Appliquer les migrations
npm run seed                  # Données initiales
npm run dev                   # Démarrer en développement
\end{lstlisting}

\subsection{Installation du frontend}
\begin{lstlisting}[language=bash,caption={Installation du frontend}]
cd web-app
npm install
cp .env.example .env          # Configurer VITE_API_URL
npm run dev                   # Démarrer en développement
npm run build                 # Build de production
\end{lstlisting}

\subsection{Installation du service IA}
\begin{lstlisting}[language=bash,caption={Installation du service IA}]
cd ai-service
pip install -r requirements.txt
python training/train_classifier.py   # Entraîner le modèle
uvicorn app.main:app --host 0.0.0.0 --port 8000
\end{lstlisting}

\subsection{Installation de l'application Flutter}
\begin{lstlisting}[language=bash,caption={Installation de l'application Flutter}]
cd worker-app
flutter pub get
# Configurer l'URL du backend dans lib/core/constants/env.dart
flutter run                   # Développement
flutter build apk --release   # Build APK de production
\end{lstlisting}


\section{Schéma de base de données}
\label{sec:schema_bdd}
Le schéma complet de la base de données est défini dans le fichier \texttt{backend/prisma/schema.prisma} (446~lignes). Il contient 14~enums et 17~modèles. Les principales relations sont :
\begin{itemize}
    \item \texttt{User} $\leftrightarrow$ \texttt{EmployeeProfile} : relation 1:1 pour les profils employés.
    \item \texttt{Room} $\leftrightarrow$ \texttt{Reservation} : relation 1:N pour les réservations par chambre.
    \item \texttt{Reservation} $\leftrightarrow$ \texttt{GuestForm} : relation 1:N pour les fiches voyageurs.
    \item \texttt{Reservation} $\leftrightarrow$ \texttt{Complaint} : relation 1:N pour les réclamations.
    \item \texttt{Complaint} $\leftrightarrow$ \texttt{InterventionLog} : relation 1:N pour les interventions.
    \item \texttt{Complaint} $\leftrightarrow$ \texttt{InternalMessage} : relation 1:N pour la messagerie interne.
    \item \texttt{Room} $\leftrightarrow$ \texttt{DailyCleaningTask} : relation 1:N pour les tâches de ménage quotidien.
    \item \texttt{User} $\leftrightarrow$ \texttt{WorkerShiftSchedule} : relation 1:N pour les plannings.
\end{itemize}


\section{Résultats détaillés de la classification IA}
\label{sec:resultats_ia}

\subsection{Rapport de classification par catégorie}

\begin{table}[H]
\centering
\caption{Rapport de classification détaillé --- Modèle LR\_combined}
\small
\begin{tabular}{|l|c|c|c|c|}
\hline
\textbf{Catégorie} & \textbf{Précision} & \textbf{Rappel} & \textbf{F1-score} & \textbf{Support} \\ \hline
COMPLAINT & 0.88 & 0.88 & 0.88 & 24 \\ \hline
HOUSEKEEPING & 0.90 & 0.90 & 0.90 & 20 \\ \hline
MAINTENANCE & 0.93 & 0.93 & 0.93 & 27 \\ \hline
OTHER & 0.68 & 0.72 & 0.70 & 18 \\ \hline
RECEPTION & 0.81 & 0.81 & 0.81 & 16 \\ \hline
RESTAURANT & 0.94 & 0.94 & 0.94 & 16 \\ \hline
\hline
\textbf{Macro avg} & 0.86 & 0.86 & 0.86 & 121 \\ \hline
\textbf{Weighted avg} & 0.86 & 0.86 & 0.86 & 121 \\ \hline
\end{tabular}
\label{tab:classification_detail}
\end{table}

\subsection{Observations}
\begin{itemize}
    \item Les catégories \texttt{MAINTENANCE}, \texttt{RESTAURANT} et \texttt{HOUSEKEEPING} obtiennent les meilleurs scores (F1 $>$ 0.90), car les réclamations associées contiennent des termes distinctifs.
    \item La catégorie \texttt{OTHER} présente les scores les plus faibles (F1 = 0.70), ce qui est attendu pour une catégorie résiduelle.
    \item La validation croisée 5-fold (CV F1 macro = 0.7702 $\pm$ 0.0450) confirme la robustesse du modèle mais indique une variabilité liée à la taille limitée du dataset.
\end{itemize}


\section{Structure du projet}
\label{sec:structure_projet}
\begin{lstlisting}[caption={Arborescence simplifiée du projet LoomStay}]
hotel-smart-concierge/
+-- backend/
|   +-- prisma/schema.prisma        # Schema BDD (446 lignes)
|   +-- src/
|   |   +-- controllers/            # 14 controllers
|   |   +-- services/               # 17 services
|   |   +-- routes/                 # 20 fichiers de routes
|   |   +-- middlewares/            # auth, role, validate, error
|   |   +-- socket/                 # Socket.IO (messagerie)
|   |   +-- utils/                  # audit, encryption, permissions
|   +-- tests/                      # 7 fichiers de tests
+-- web-app/
|   +-- src/
|   |   +-- pages/staff/            # 3 dashboards + login
|   |   +-- pages/public/           # 8 pages client
|   |   +-- components/             # Composants réutilisables
|   |   +-- routes/                 # Configuration React Router
+-- ai-service/
|   +-- app/
|   |   +-- main.py                 # FastAPI endpoints
|   |   +-- services/               # classifier, translation, language
|   +-- training/                   # Script d'entraînement
|   +-- models/                     # Modèle sérialisé (.joblib)
|   +-- tests/                      # Tests Pytest
+-- worker-app/
|   +-- lib/
|   |   +-- features/               # auth, home, tasks, splash
|   |   +-- core/                   # API client, router, constants
|   |   +-- shared/                 # Widgets réutilisables
+-- dataset/                        # Dataset de réclamations (604 ex.)
\end{lstlisting}
